\documentclass{otbook}
% Package loading: this class loads its companion bundle
% (otnotation, otmath, otbib, otindex, ...) for you.
% For a lean build, use \documentclass[minimal]{otbook} and load only
% what you need, e.g. \usepackage{otphysics}. The semantic boxes
% (definitionbox, formulabox, ...) are available either way.

\typeout{OTBOOK-STARTER:PERSONAL-BOOK}

\title{Personal Book Title}
\booksubtitle{A working subtitle for the project}
\author{Author Name}
\bookedition{Draft edition}
\date{\today}

\begin{document}

\frontmatter
\hypersetup{pageanchor=false}
\makeotbooktitle%
% \makeotbookdedication{For the people who made this book possible.}
\hypersetup{pageanchor=true}
\tableofcontents

\mainmatter%

\chapter{Foundations}

\otbookepigraph{Add a short quotation that frames the chapter.}{Attribution}

\begin{chapteroverview}
  Explain the question, boundary, or thread that this chapter will develop.
\end{chapteroverview}

\begin{chaptergoals}
  \item State the central idea of the chapter.
  \item Record the formulae or methods that should be easy to find later.
  \item Add one worked example and one personal reflection.
\end{chaptergoals}

\section{Main Idea}

\begin{definitionbox}[Core Definition]
  Write the definition, principle, or concept that anchors the chapter.
\end{definitionbox}

\begin{intuitionbox}[How To Remember It]
  Add the short explanation that would help you recover the idea after time away.
\end{intuitionbox}

\section{Formula Pattern}

\begin{formulabox}[A Useful Relationship]
  \[
    y=mx+c
  \]
\end{formulabox}

\begin{examplebox}[Worked Example]
  If \(m=2\), \(c=1\), and \(x=3\), then \(y=2(3)+1=7\).
\end{examplebox}

\begin{reflectionbox}
  Add the personal note, caveat, or connection that makes this chapter yours.
\end{reflectionbox}

\backmatter

\chapter{Reference Notes}

\begin{summarybox}
  Keep formula references, reading notes, or revision checklists here.
\end{summarybox}

\end{document}
